% TOP_PROBLEM: {"id": 30007229, "title": "The area, charge and angular-momentum bound with controlled electromagnetic tails", "category_id": 16, "category": {"id": 16, "name": "physics", "display_name": "Mathematical Physics", "description": "Problems at the intersection of mathematics and physics.", "slug": "physics", "order_index": 16, "created_at": "2026-07-31T15:26:25.671Z"}, "status": "open"}
% FIRST_POSTED: 2026-10-06
% ATTEMPT_STATUS: solved
% OPENAI_CLAIM: {"schema_version":1,"source_repo":"https://github.com/openai/math","source_commit":"adc7f1241b42e322a6451854ab7e4b4c146bf78a","release_date":"2026-10-06","match":"full","resolution":"proved","independently_reviewed":false,"second_pass":"agreed","families":[{"family":"260","title":"Spacetime Penrose inequalities: enclosing area, charge, rotation, and anti-de Sitter extensions","subject":"Mathematical physics","manuscripts":[{"title":"The Kerr–Newman Penrose Inequality for Axisymmetric Electrovacuum Exteriors","pdf_path":"preprints/The-Kerr-Newman-Penrose-Inequality-for-Axisymmetric-Electrovacuum-Exteriors-October-5-2026/paper.pdf","readme_path":"preprints/The-Kerr-Newman-Penrose-Inequality-for-Axisymmetric-Electrovacuum-Exteriors-October-5-2026/README.md","date":"2026-10-05","theorem_ref":"preprints/The-Kerr-Newman-Penrose-Inequality-for-Axisymmetric-Electrovacuum-Exteriors-October-5-2026/build/main.tex:293 (m:main)"}],"lean":{"doc_path":"lean/docs/260.md","comparators":[],"declarations":[],"covers_main_theorem":false},"reasoning_trace":null}]}
\documentclass[11pt]{article}
\usepackage[margin=25mm]{geometry}
\usepackage{fontspec}
\setmainfont{Latin Modern Roman}
\usepackage{amsmath,amssymb,mathtools}
\usepackage{unicode-math}
\setmathfont{Latin Modern Math}
\usepackage{enumitem}
\usepackage{xurl,hyperref,bookmark}
\hypersetup{colorlinks=true,urlcolor=blue}
\setcounter{secnumdepth}{0}
\setlength{\parindent}{0pt}
\setlength{\parskip}{5pt}
\setlength{\emergencystretch}{3em}
\begin{document}
\section*{The area, charge and angular-momentum bound with controlled electromagnetic tails}

\subsection{OpenAI's claim}
OpenAI presents a proof in the selected manuscript. For the designated axisymmetric electrovacuum class, the claimed bound is $m^2\geq A/(16\pi)+Q^2/2+\pi(Q^4+4J^2)/A$ with conserved total angular momentum.

\subsection{Outline of the argument}
Axial variables are compared with a Kerr--Newman family at prescribed conformal capacity. Riemannian scale comparison and electromagnetic angular-momentum bookkeeping give the mass bound; saturation on the strict physical-area branch reconstructs an admissible dyonic spacetime slice.

\subsection{Scope relative to this problem}
For the designated axisymmetric electrovacuum class, the claimed bound is $m^2\geq A/(16\pi)+Q^2/2+\pi(Q^4+4J^2)/A$ with conserved total angular momentum. The question retains these same parameters and conventions. Within these assumptions, this is a claimed affirmative answer to the stated question.

\subsection{Formal verification}
For family 260, the only selected family entry is the CKS end-replacement construction and Schwarzschild equality examples. A connected-boundary inequality is conditional on the flat Penrose input; the general Bondi or spacetime inequalities are not formalized there. No Comparator is assigned to the selected manuscript here; the family page selects other papers or supporting statements. This site has not executed the Lean code. The scope descriptions do not certify the remaining manuscript arguments.

\subsection{Status caveat}
OpenAI attributes these manuscripts to work by an internal model. Its release cautions that arguments without formalization can contain mistakes. This catalogue has not independently checked the proofs; the labels report the mathematical claims attributed to OpenAI.

\subsection{Sources}
\begin{itemize}
\item OpenAI, \emph{The Kerr–Newman Penrose Inequality for Axisymmetric Electrovacuum Exteriors}, 2026-10-05; citation key \texttt{OAI:The-Kerr-Newman-Penrose-Inequality-for-Axisymmetric-Electrovacuum-Exteriors-October-5-2026}. \href{https://github.com/openai/math/blob/adc7f1241b42e322a6451854ab7e4b4c146bf78a/preprints/The-Kerr-Newman-Penrose-Inequality-for-Axisymmetric-Electrovacuum-Exteriors-October-5-2026/paper.pdf}{PDF (GitHub)}; \href{https://github.com/openai/math/raw/adc7f1241b42e322a6451854ab7e4b4c146bf78a/preprints/The-Kerr-Newman-Penrose-Inequality-for-Axisymmetric-Electrovacuum-Exteriors-October-5-2026/paper.pdf}{PDF (direct)}; \href{https://github.com/openai/math/blob/adc7f1241b42e322a6451854ab7e4b4c146bf78a/preprints/The-Kerr-Newman-Penrose-Inequality-for-Axisymmetric-Electrovacuum-Exteriors-October-5-2026/README.md}{Manuscript page and citation}; \href{https://github.com/openai/math/blob/main/preprints/The-Kerr-Newman-Penrose-Inequality-for-Axisymmetric-Electrovacuum-Exteriors-October-5-2026/paper.pdf}{Latest version}.
\item Family 260: \href{https://github.com/openai/math/blob/adc7f1241b42e322a6451854ab7e4b4c146bf78a/lean/docs/260.md}{Lean scope}.
\item \href{https://github.com/openai/math/blob/adc7f1241b42e322a6451854ab7e4b4c146bf78a/CONTENTS.md}{OpenAI catalogue entry}.
\end{itemize}
\end{document}
